\section{Experimental protocol and comparison scope}
\label{sec:protocol}
\subsection{Development data}
The current synthetic training pool contains 57 DAVIS training videos,
excluding the three development sources \emph{bear}, \emph{lucia}, and
\emph{walking}. The synthetic development pilot uses 64 frames per source,
at 480-pixel height, with two independently sampled noise conditions.
There are three source videos and six conditions, not six independent videos.
For a target RGB value $y$ in $[0,1]$, the degradation proxy is
\[
 v=0.35y^{2.2},\qquad
 x=Q_8\!\left(\operatorname{clip}
 \left[\operatorname{Poisson}(Pv)/P+\eta,0,1\right]\right),
 \quad \eta\sim\mathcal N(0,0.008^2),
\]
with $P\in\{20,80\}$. Here $Q_8$ denotes conversion to the eight-bit
observation. The present pilot adds no extra JPEG encoding, although the
source frames are JPEG images. This is a controlled synthetic proxy, not
a calibrated model of arbitrary RAW sensors or phone ISP video.

The fixed mixed-data study replaces half of each training batch with
native crops from 12 verified paired ND-filter video sources in the
project's OLiVES snapshot. Its four already-used paired-real development
sources are videos 15, 47, 91, and 121, each evaluated on 16 frames with a
$288\!\times\!512$ crop. Training and development source IDs are separated.
Consecutive registered frame indices are available; independent sensor
timestamps have not been verified. These daytime ND pairs must not be
presented as natural nighttime validation.

Three original phone clips have additionally been inspected without
synthetic degradation. They have no synchronized reference, so they support
behavioural examples rather than reference PSNR or a perceptual flicker
ranking. Independent final-test evaluation remains outstanding.

\subsection{Quality and temporal measurements}
Predictions are evaluated as floating-point RGB values after clipping to
$[0,255]$, before display rounding or video encoding. Sequence PSNR is the
mean of full-frame RGB frame PSNR values,
\[
 \operatorname{PSNR}_{\rm seq}
 =\frac1T\sum_t 10\log_{10}
 \frac{255^2}{\operatorname{MSE}_{\rm RGB}(z_t,y_t)}.
\]
RGB SSIM uses an $11\!\times\!11$ Gaussian window with standard deviation
1.5 and the valid region obtained by excluding five boundary pixels.
Edge MAE is the full-frame mean absolute error of horizontal and vertical
luminance Sobel gradients, with the fixed gradient normalization used by
the evaluation implementation.

For a fixed estimated reference flow and validity mask, temporal error is
\[
 E_{\rm temp}=\frac{1}{|\mathcal T|}\sum_{t\in\mathcal T}
 \frac{\sum_{p,c}m_t^{\rm ref}(p)
 \left|[z_t-W_t^{\rm ref}z_{t-1}]_c(p)
       -[y_t-W_t^{\rm ref}y_{t-1}]_c(p)\right|}
 {3\sum_p m_t^{\rm ref}(p)}.
\]
Here $\mathcal T$ contains frame pairs with nonempty valid masks. The
sequence summary first averages each valid frame pair, then averages
the available pairs. The reference mask excludes an
eight-pixel boundary. Reference flow is estimated from targets and held
fixed across compared methods; it is not ground-truth motion.
This metric is not a sufficient perceptual definition of flicker.
Evaluation uses fixed reference fields and absolute RGB error. Training
uses half-scale estimated fields, differentiable continuous bilinear
gathering, and squared error. The two are not numerically identical.

Domains and training seeds are reported separately before aggregation.
The present two-seed pilot does not establish statistical significance.
We retain per-condition failures alongside mean improvements.

\subsection{Strong same-weight references}
For each spatial anchor, let $r_t$ denote the raw prediction. The finite
reference uses the previous raw output, while recursive fusion uses the
previous fused output:
\[
 z_t^{\rm fin}=\tfrac12 r_t+\tfrac12 W_t r_{t-1},
 \qquad
 z_t^{\rm rec}=\tfrac12 r_t+\tfrac12 W_t z_{t-1}^{\rm rec}.
\]
The formulas apply only on accepted correspondences; invalid points and the
first frame use the current raw output. Both estimate motion from their own
adjacent raw predictions with the fixed half-resolution forward--backward
estimator. Recursive fusion therefore also keeps the previous raw output.
Neither reference requires a second enhancement CNN or a GT reset.

Published FastDVDnet weights~\cite{fastdvdnet} have been evaluated as
fixed denoiser transfer, not as a retrained low-light video method.
Its five-frame window uses two future frames. The documented pretraining
source includes the three current DAVIS development IDs, so that transfer
cannot establish unseen generalization. Same-domain retraining and broader
published LLVE comparisons remain necessary.

\subsection{Efficiency}
Convolution MACs, pointwise work, history storage, and full pipeline time
are distinct measurements. The original predictor has approximately
7.938 convolution GMAC at $1080\mathrm p$. This count omits non-convolution
work. Online history counts use actual tensor storage rather than only
logical channel slices. CPU timing is measured with two Torch threads and
one OpenCV thread, and its workload and reference weights are stated
explicitly. No handset latency is inferred from these measurements.
